% !TEX root = ../main.tex
\chapter{Fondamenti Teorici}
\label{chap:logical_foundations}

\section{Il \texorpdfstring{$\lambda$}{Lambda}-calcolo}
Il $\lambda$-calcolo è un formalismo computazionale, ovvero un insieme di regole logico-matematiche 
precise utilizzate per definire le funzioni calcolabili. Il concetto stesso di funzione calcolabile è 
definita su questi formalismi secondo la tesi di Church-Turing. 

\subsection{Grammatica del \texorpdfstring{$\lambda$}{Lambda}-calcolo}
La grammatica BNF del $\lambda$-calcolo, definita in \eqref{eq:lambda_grammar}, si basa su tre costrutti fondamentali:
le varibili, l'astarzione e l'applicazione. La grammatica definisce solo il sistema di scrittura dei termini, 
essa non assegna alcuna semantica e nessuna relazione tra gli stessi ad esclusione di quelle sintattiche.
\begin{equation}
    M, N ::= x \mid (\lambda x. M) \mid (M N)
    \label{eq:lambda_grammar}
\end{equation}
In ordine, i costrutti prendono il nome di variabile, astrazione e applicazione; la loro arieta è rispettivamete atomina, unaria e binaria.
L'arieta all'interno della grammatica deternima il numero di elementi richiesti per la loro formazione dei termini. È inoltre opportuno osservare 
che la definizione di queste proprietà appartiene almeno per ora al metalinguaggio, ovvero il linguaggio logico-matematico e non al $\lambda$-calcolo.
    
\subsection{Relazioni tra termini}
Le relaioni tra termini del $\lambda$-calcolo sono cio che traforma una insieme di puramente sintattico in un effettivo formalismo computazionale. 
Alcune di esse sono necessarie 

\subsubsection{Definizioni preliminari}
% TODO: Variabili libere e legate, Sostituzione, Sottotermini

\subsubsection{$\alpha$-equivalenza}
la $\alpha$-equivalenza stabilisce che il nome delle variaibli legate è irrilevante ai fini del significato del ternine. Due termini sono considerati 
identici se l'uno puo essere ottenuto dal'altro tramite la ridenominazione coerente delle proprie variabili legate. Questa relazione è fondametare per 
evitare il fenomeno del varible capture durante le operazione di sostituzione, garantendo che i nomi scelti arbitrariamente non interferiscano con 
le varibili libere.
\begin{equation}
    \lambda x. M \equiv_\alpha \lambda y. M[y/x] \quad \text{se } y \notin FV(M)
    \label{eq:alpha_conversion}
\end{equation}

\subsubsection{$\beta$-riduzione}
\label{sec:beta_reduction}
La $\beta$-riduzione stabilisce il motore del $\lambda$-calcolo e formalizza il cncetto di applicazione funzionale. Consiste nel sostituire l'argomento N a ogni
occorenza libera della variabile x allinterno del corpo della funzione M.
\begin{equation}
    (\lambda x. M) N \to_\beta M[x := N]
    \label{eq:beta_reduction}
\end{equation}
Un ternine che presenta questa struttura si dice \emph{redex} (reducible expression). Il calcolo consiste dell'applicazione ripetura di questa regola, 
quando un termine non contiene piu redex si dice che è in \emph{forma normale}. La forma normale è il risultato concreto della computazione.

\subsubsection{$\eta$-conversione}
La $\eta$-conversione ($\to_\eta$) introduce il principio di estensionalità, secondo il quale due funzioni sono uguali se, 
applicate allo stesso argomento, producono il medesimo risultato.
\begin{equation}
    \lambda x. (M x) \to_\eta M \quad \text{se } x \notin FV(M)
    \label{eq:eta_conversion}
\end{equation}
Questa operazione permette di semplificare espressioni ridondanti, eliminando astrazioni che si limitano a "passare" 
l'argomento a un altro termine.

\subsection{Problemi di divergenza}
Nel $\lambda$-calcolo, come esposto nella sezione \ref{sec:beta_reduction}, la computazione deriva dalla riduzione dei termini. Quando questa non è possibile allora
il termine è in forma normale, ovvero il risultato della computazione. Tuttavia è importante osservare che non tutti i termini ammetono una forma normale, in alcuni 
casi la riduzione può proseguire indefinitivamente. 


\subsubsection{Definizione di divergenza}
L'assenza di una forma normale è ciò che definisce formalmente la \emph{divergenza}. 
In un sistema Turing-completo come il $\lambda$-calcolo non tipato, la divergenza è 
l'equivalente funzionale di un ciclo infinito (loop) in un linguaggio di programmazione imperativo.

\subsubsection{Il termine $\Omega$}
L'esempio piu semplice di termine divergente è il combinatore $\Omega$. Esso viene costruito a partire dal termine di auto-applicazione $\omega$ 
(talvolta indicato come $\delta$).
\begin{equation}
    \omega = \lambda x. (x x)
\end{equation}
Dalla regola di $\beta$-riduzione appare evidente che l'auto-applicazione del termine $\omega$ sia divergente, la contrazione del redex non conduce a una forma normale ma 
riporta al termine iniziale.

\section{Sistemi di Tipi e Normalizzazione}

\subsection{Il \texorpdfstring{$\lambda$}{Lambda}-calcolo semplicemente tipato (\texorpdfstring{$\lambda^{\to}$}{Lambda-to})}
Il $\lambda$-calcolo semplicemente tipato ($\lambda^{\to}$), come suggerisce il nome, si distingue dal sistema puro per l'introduzione dei tipi.
In questo sistema di riscrittura i tipi sono della annotazioni sulla varibile vincolata della astrazione

\subsubsection{Grammatica del \texorpdfstring{$\lambda$}{Lambda}-calcolo semplicemente tipato}
La grammatica BNF del $\lambda$-calcolo semplicemente tipato ($\lambda^{\to}$), estende il calcolo puro tramite
la categoria sintattica dedicata ai tipi \eqref{eq:simple_type_grammar}, legando le categorie sintattice tramite 
il costrutto dell'astrazione \eqref{eq:simple_typed_lambda_grammar}.
\begin{equation}
    \tau ::= \iota \mid \tau_1 \to \tau_2
    \label{eq:simple_type_grammar}
\end{equation}
\begin{equation}
    e ::= x \mid \lambda x : \tau . e \mid e \, e
    \label{eq:simple_typed_lambda_grammar}
\end{equation}
Si osservi come la sintassi dell'astrazione includa il tipo, cio garantiche che ogni variabile vincolata sia univovamete tipata 
all'interno del proprio scope. Inoltre è opportuno ribadire che la grammatica non fornisce semantica.

\subsubsection{Il giudizio di tipizzazione o type checking}
Il giudizio di tipizzazione o \emph{type checking} ($\Gamma \vdash e : \tau$) è il processo d'\emph{inferenza} mediante il quale si determina se un termine appertiene a un determinato tipo sotto un insieme di assunzioni.
Tale processo opera in modo induttivo sulla grammatica dei termini, visitando l'intera struttura verificandone le componenti. Le applicazioni ricorsive delle regole per verificare il tipo finale,
garantendo la correttezza formale dell'intera computazione.
\begin{equation}
    \text{type}(\Gamma, e) = 
    \begin{cases} 
        \tau & \text{if } e = x \text{ e } (x : \tau) \in \Gamma \\
        \sigma \to \text{type}(\Gamma, x:\sigma, e') & \text{if } e = \lambda x : \sigma . e' \\
        \tau & \text{if } 
            \begin{cases} 
            e = e_1 e_2 \\
            \text{type}(\Gamma, e_1) = \sigma \to \tau \\
            \text{type}(\Gamma, e_2) = \sigma 
            \end{cases} \\
        \perp & \text{otherwise}
    \end{cases}
    \label{eq:simple_typed_lambda_type_cheking_algorithm}
\end{equation}
Il contesto di tipizzazione, denotato con $\Gamma$, è un insieme finito di assunzioni di tipo nella forma $x : \sigma$.
L'estensione del contesto, indicata con $\Gamma, x : \sigma$, rappresenta l'aggiunta di una nuova assunzione, con la clausola 
che se $x$ era già presente in $\Gamma$, la nuova assegnazione sovrascrive la precedente (\emph{meccanismo di shadowing}). Il contesto
può essere modellizato come una funzione parziale. Il simbolo $\perp$ (leggasi bottom) invece rapresenta il fallimento dell'algoritmo
di verifica, questo in seguito rapresentera una specifica condizione sul termine, limitiamoci a dire che è mal tipato (\emph{ill-typed}).

\subsubsection{L'inferenza di tipo o type inference}
L'inferenza di tipo o \emph{type inference}  è il prcesso mediante il quale si ricostruisce automanticamente il tipo piu generale di un termine,
utilizando al struttura sintattica e le relazioni tra componenti.

\begin{equation}
    \text{gen}(\Gamma, e) = 
    \begin{cases} 
        (\tau, \emptyset) & \text{if } e = x \land (x : \tau) \in \Gamma \\
        (\alpha \to \tau, \mathcal{C}) & \text{if } e = \lambda x. e' \land (\tau, \mathcal{C}) = \text{gen}(\Gamma, x : \alpha, e') \\
        (\beta, \mathcal{C}_{all}) & \text{if } e = e_1 e_2 \text{ where:} \\ & 
        \begin{cases} 
            (\tau_1, \mathcal{C}_1) = \text{gen}(\Gamma, e_1) \\
            (\tau_2, \mathcal{C}_2) = \text{gen}(\Gamma, e_2) \\
            \mathcal{C}_{all} = \mathcal{C}_1 \cup \mathcal{C}_2 \cup \{\tau_1 = \tau_2 \to \beta\}
        \end{cases}
    \end{cases}
    \label{eq:constraint_generation}
\end{equation}
Il risultato del processo d'inferenza di tipo è un insieme di equazioni $\mathcal{C}$ che rapresentano i vincoli tra tipi delle componenti del termine.
Se il sistema di equazioni dei vicoli di tipi ha soluzione allora è possibile dimostrare che il risoltato è il tipo piu generale che rispetta tali vincoli 
(\emph{MGU}). L'insieme di equasioni $\mathcal{C}$ viene anche chiamato \emph{probela di unificazione del primo ordine}.
\begin{equation}
    \mathcal{U}(\mathcal{C}) = 
    \begin{cases} 
        id & \text{if } \mathcal{C} = \emptyset \\
        \mathcal{U}(\mathcal{C}') & \text{if } \mathcal{C} = \{\tau = \tau\} \cup \mathcal{C}' \\
        \mathcal{U}(\mathcal{C}'[\alpha \mapsto \tau]) \circ [\alpha \mapsto \tau] & \text{if } \mathcal{C} = \{\alpha = \tau\} \cup \mathcal{C}' \text{ e } \alpha \notin FV(\tau) \\
        \mathcal{U}(\mathcal{C}'[\alpha \mapsto \tau]) \circ [\alpha \mapsto \tau] & \text{if } \mathcal{C} = \{\tau = \alpha\} \cup \mathcal{C}' \text{ e } \alpha \notin FV(\tau) \\
        \mathcal{U}(\{\tau_1 = \sigma_1, \tau_2 = \sigma_2\} \cup \mathcal{C}') & \text{if }    \mathcal{C} = \{\tau_1 \to \tau_2 = \sigma_1 \to \sigma_2\} \cup \mathcal{C}' \\
        \perp & \text{otherwise (failed)}
    \end{cases}
    \label{eq:unification_algorithm}
\end{equation}
La risoluzione del sistema di vincoli poggia su principi logici intuitivi ma rigorosi: 
l'impossibilità di far coincidere tipi atomici distinti (assenza di collisioni tra costruttori) e 
il divieto di formulare definizioni ricorsive, per cui un tipo non può essere definito in funzione di se stesso (occurs check).
\begin{table}[h]
    \centering \small
    \begin{tabular}{lp{5cm}l}
        \hline
        \textbf{Errore} & \textbf{Descrizione}                   & \textbf{Esempio di Vincolo}  \\ \hline
        Clash Error     & Unificazione di costruttori diversi    & $Bool = \sigma \to \tau$     \\
        Occurs Check    & Definizione di tipi ricorsivi infiniti & $\alpha = \alpha \to \beta$  \\
        Type Mismatch   & Incongruenza tra tipi base differenti  & $Nat = Bool$                 \\ \hline
    \end{tabular}
    \caption{Fallimenti dell'algoritmo di unificazione.}
    \label{tab:unification_failures}
\end{table}

\subsubsection{Normalizzazione forte (Teorema)}


% \section{L'Isomorfismo di Curry-Howard}